\documentclass[11pt,a4paper]{coverletter}

\senderblock
  {Radheem Bin Razi}
  {Distributed Systems \& Cloud-Native Engineer}
  {Ilmenau, Germany \textbar\ \href{mailto:sheikh.radheem@gmail.com}{sheikh.radheem@gmail.com} \textbar\ \href{https://radheem.github.io/my-cv}{portfolio}}

\begin{document}

%% ===== ENGLISH =====
\recipient{Moss}{Hiring Team}{}

\opening{Dear Hiring Team,}

Moss is making a deliberate bet on AI tooling and hiring someone to turn that bet into measurable outcomes — and that is exactly the kind of open-ended, build-it-from-first-principles work I gravitate toward. I am a senior software engineer who has spent the last few years building distributed services for fintech and, more recently, building the agentic tooling and guardrails that let AI sit safely alongside real engineering work. My strengths are the three things this role needs at once: hands-on platform engineering, practical experience with MCP and LLM integrations, and the judgement to weigh productivity gains against the safety and compliance constraints a regulated business cares about.

A large part of my recent work has been about exposing systems to AI agents responsibly. On a distributed Go platform, I built MCP tooling that lets LLM agents drive the system through well-defined interfaces rather than ad-hoc access — the same kind of safe, governed surface area Moss wants for IDE integrations, MCP servers, and model integrations like Claude and Codex. The deeper lesson there is the one your guardrails responsibility is really about: AI is only as useful as the boundaries around it are trustworthy.

I have also lived the productivity-versus-safety trade-off in a smaller, sharper form. I built a Python tool where a pure, unit-tested component makes all the consequential decisions and the LLM is confined to writing prose, so it cannot fabricate facts. It uses git and GitHub Actions as its system of record, with no secrets in CI and an encrypted gate over anything sensitive. That instinct — let AI accelerate the work while a deterministic, observable layer keeps it honest — is what I would bring to Moss's approved-model catalogue, prompt logging, and usage observability. Earlier, as a design authority on a high-throughput crypto exchange, I learned to identify what needs solving across teams and drive the fix with their buy-in, rather than waiting for a ticket.

I would welcome the chance to do this work at Moss. I am based in Ilmenau, Germany, available full-time now, and can relocate within Germany in roughly two to three weeks. I am Blue Card-ready, so a qualifying offer needs only the contract — no sponsorship overhead. Thank you for considering my application; I would be glad to talk through how I can help make AI adoption at Moss both faster and safer.

\closing{Sincerely,}{Radheem Bin Razi}

\clearpage
\selectlanguage{ngerman}

%% ===== DEUTSCH =====
\recipient{Moss}{Personalabteilung}{}

\opening{Sehr geehrtes Hiring-Team,}

Moss setzt bewusst auf KI-Tooling und sucht jemanden, der diese Wette in messbare Ergebnisse verwandelt — und genau das ist die Art von offener, von Grund auf zu konzipierender Arbeit, zu der es mich hinzieht. Ich bin ein Senior Software Engineer, der in den letzten Jahren verteilte Dienste für den Fintech-Bereich aufgebaut hat und zuletzt das agentenbasierte Tooling sowie die Leitplanken entwickelt hat, die es KI ermöglichen, sicher neben echter Engineering-Arbeit zu bestehen. Meine Stärken sind die drei Dinge, die diese Rolle gleichzeitig erfordert: praxisnahes Platform Engineering, fundierte Erfahrung mit MCP und LLM-Integrationen sowie das Urteilsvermögen, Produktivitätsgewinne gegen die Sicherheits- und Compliance-Anforderungen abzuwägen, die einem regulierten Unternehmen wichtig sind.

Ein großer Teil meiner jüngsten Arbeit drehte sich darum, Systeme verantwortungsvoll für KI-Agenten zugänglich zu machen. Auf einer verteilten Go-Plattform habe ich MCP-Tooling entwickelt, das LLM-Agenten erlaubt, das System über klar definierte Schnittstellen statt über Ad-hoc-Zugriff zu steuern — dieselbe Art sicherer, geregelter Angriffsfläche, die Moss für IDE-Integrationen, MCP-Server und Modellintegrationen wie Claude und Codex anstrebt. Die tiefere Lektion daraus ist genau das, worum es bei Ihrer Verantwortung für die Leitplanken wirklich geht: KI ist nur so nützlich, wie die Grenzen um sie herum vertrauenswürdig sind.

Ich habe den Kompromiss zwischen Produktivität und Sicherheit auch in einer kleineren, schärferen Form erlebt. Ich habe ein Python-Tool entwickelt, bei dem eine reine, unit-getestete Komponente alle folgenreichen Entscheidungen trifft und das LLM auf das Schreiben von Prosa beschränkt ist, sodass es keine Fakten erfinden kann. Es nutzt git und GitHub Actions als System of Record, ohne Secrets in der CI und mit einem verschlüsselten Gate über allem Sensiblen. Dieser Instinkt — KI die Arbeit beschleunigen zu lassen, während eine deterministische, beobachtbare Schicht sie ehrlich hält — ist das, was ich in Moss' Katalog freigegebener Modelle, Prompt-Logging und Usage-Observability einbringen würde. Zuvor habe ich als Design-Authority an einer Krypto-Börse mit hohem Durchsatz gelernt, teamübergreifend zu erkennen, was gelöst werden muss, und die Lösung mit deren Zustimmung voranzutreiben, anstatt auf ein Ticket zu warten.

Ich würde die Gelegenheit begrüßen, diese Arbeit bei Moss zu leisten. Ich bin in Ilmenau, Deutschland, ansässig, ab sofort in Vollzeit verfügbar und kann innerhalb Deutschlands in etwa zwei bis drei Wochen umziehen. Ich bin Blue-Card-ready, sodass ein qualifizierendes Angebot nur den Vertrag benötigt — ohne Sponsoring-Aufwand. Vielen Dank, dass Sie meine Bewerbung in Betracht ziehen; ich würde mich freuen, zu besprechen, wie ich dazu beitragen kann, die KI-Einführung bei Moss sowohl schneller als auch sicherer zu gestalten.

\closing{Mit freundlichen Grüßen,}{Radheem Bin Razi}

\end{document}
